%% kapa: exam=1 round=24 year=2013 period=1 form=A subject=economics subject_ko=경제원론 seq=10 number=50
%% source: Q-Net 감정평가사 기출문제 (cst00309) artlSeq=5073802 file=extracted/1008/5073802/2013년 제24회 감정평가사 1차시험 기출문제/1교시(A형).hwp
%% conversion: kordoc 4.24.51 -> tools/gen1.py (auto); review=proofread-sonnet; flags=none
\begin{question}{50}
A기업의 생산함수는 $Q= \sqrt { KL }$, 자본($K$)의 가격은 $r$, 노동($L$)의 가격은 $w$이다. 생산량이 $Q_{ 0 }$로 주어졌을 때, 비용이 극소화되도록 자본과 노동의 투입량을 결정하고자 한다. 이에 관한 설명으로 옳지 않은 것은? (단, $Q_{ 0 }$>0이다.)
\begin{choices}
  \item 생산함수는 자본과 노동에 대해 1차 동차함수이다.
  \item 최적 상태에서 노동 1단위당 자본 투입량은 $\frac{ r } { w }$이다.
  \item 최적 상태에서 요소간의 대체탄력성은 1이다.
  \item 최적 상태에서 노동과 자본의 투입량은 $w$, $r$, $Q_{ 0 }$의 함수이다.
  \item 최적 상태에서 총비용 중 노동비용이 차지하는 비중은 일정하다.
\end{choices}
\end{question}
\answer{2}
